\answer
\begin{enumerate}

%% 2章 問1（2.1節）
\item[2.1節 問1]
導通損失は式\ref{eq2.1}より
\begin{equation*}
P_{on} = R_{on}I^2 = 2\times10^{-3}\times50^2 = 5\,\mathrm{W}
\end{equation*}
オフ状態の損失は$600\times10^{-6} = 0.6$\,mWで，導通損失の1万分の1程度である。
よく設計された素子ではオフの損失は無視でき，
損失の主役は導通損失とスイッチング損失である。

%% 2章 問2（2.2節）
\item[2.2節 問2]
(a) ダイオード：制御端子なし。BJT・MOSFET・IGBT：自己消弧形。
サイリスタ：他励形（オンのみ制御できる）。
(b) BJT・サイリスタ：電流駆動。MOSFET・IGBT：電圧駆動。
ダイオードは駆動そのものがない。

%% 2章 問3（2.3節）
\item[2.3節 問3]
分子$e v_D$の単位は$\mathrm{C\times V} = \mathrm{C\cdot V} = \mathrm{J}$。
分母$kT$の単位は$\mathrm{J/K\times K} = \mathrm{J}$。
よって$\dfrac{e v_D}{kT}$はJ/Jとなり無次元である。
指数関数の引数は必ず無次元でなければならないから，
この検算は式の形の正しさの確認にもなっている。

%% 2章 問4（2.4節）
\item[2.4節 問4]
コレクタ電流のもとは，エミッタからベースへ注入される電子である。
この注入が続くには，エミッタ-ベース接合の順バイアス，
すなわちベース領域に正孔が十分あることが必要である。
ベースでは電子の一部が正孔と再結合して正孔を消費するから，
ベース電流はこの消えた正孔を外から補充している。
$i_B = 0$にすると補充が絶たれて正孔が減り，
順バイアス状態が保てなくなって電子の注入が止まる。
その結果コレクタ電流も止まる。

%% 2章 問5（2.5節）
\item[2.5節 問5]
ゲートは酸化膜（絶縁体）の上にあるので，
定常状態ではゲートに電流が流れず，駆動電力はほぼゼロである。
一方，ゲート電極と半導体は酸化膜をはさんだコンデンサ（ゲート容量）をなすから，
オン・オフを切り替える瞬間には，
このコンデンサを充電・放電するための電流がゲートに流れる。
したがって駆動電力はスイッチング周波数に比例して増える。

%% 2章 問6（2.6節）
\item[2.6節 問6]
IGBTの主電流は，裏面のp形コレクタ層とn$^-$ドリフト層の間のpn接合を
順方向に通り抜ける。
順バイアスのpn接合には約$0.7$\,Vの順方向電圧が必ず立つから（\ref{sec2.3}節），
IGBTのオン電圧にはこの分の下駄が乗る。
MOSFETでは電子がソースから反転層を通ってドレインまで
pn接合の障壁を一度も越えずに流れるため，この下駄がない。

%% 2章 問7（2.7節）
\item[2.7節 問7]
サイリスタの4層は，npnトランジスタとpnpトランジスタが
中央の2層を共有した構造とみなせる。
オンの状態では，npnのコレクタ電流がpnpのベース電流となり，
pnpのコレクタ電流がnpnのベース電流となって，
2つのトランジスタが互いの駆動電流を供給し合っている。
このループが自立しているため，
外部からのゲート電流を切っても導通は保たれる（ラッチ動作）。
オフにするには，主回路の電流そのものを外部で減らして
ループを断ち切るしかない。

\end{enumerate}
